\section{Spieler programmieren}

\begin{minipage}[t]{0.56\textwidth}
\vspace{-11cm}
Probiere nach jedem Schritt aus, was passiert, wenn du dein Spiel startest!
\begin{enumerate}
    \item Programmiere die Klasse Spieler wie im Klassendiagramm rechts vorgegeben. Attribute erstellst du über \emph{Variable $\rightarrow$ Create variable} (und anschließend Eingabe des Namens.\\ 
    $\rightarrow$Programmiere getX() und getY() als Standard-Getter, lasse run() erstmal leer.
    \item Setze in run(), die Koordinaten auf feste Werte und finde heraus, wie groß das Spielfeld ist. Den Wert setzt man so: \includegraphics[height= 15pt]{Aufgaben/BlocksAndJava/img/setX123.png} 
    \item Notiere den Programmcode von getX() und getY() auf der Rückseite.
    \item Das Spielfeld führt die Methode run() seiner Spieler-Objekte ca. 25x pro Sekunde aus. Finde heraus, wie du ein Spieler-Objekt auf dem Spielfeld bewegen kannst.
    \item Wie kann man die Bewegungsgeschwindigkeit des Spielers festlegen? Probiere deine Idee im Programm aus und ergänze das Klassendiagramm und den Programmcode auf der Rückseite.
    \item Wie kann man zu Beginn des Spiels die Startposition des Spielers festlegen? Probiere auch das wieder aus und ergänze das Klassendiagramm.
    \item Das Spielfeld kann außerdem den Rückgabewert einer Methode getText() von seinen Spielerobjekten anzeigen. Überlege dir eine solche Methode und ergänze sie auf der Rückseite und im Diagramm.
    \item Das Spielfeld überprüft, ob der Spieler mit einer Münze zusammengestoßen ist und ruft, falls ja, die Methode crashMuenze() des Spieler-Objekts auf. Ausprobieren$\rightarrow$Klassendiagramm$\rightarrow$Code :)
\end{enumerate}
\end{minipage}
\hfill
\begin{minipage}[t]{0.42\textwidth}

\includegraphics[width=0.9\linewidth]{Aufgaben/BlocksAndJava/img/57_Klassendiagramm.png}
\end{minipage}

\large
\begin{lstlisting}
public class Spieler {
    private double x;
    private double y;

    (*@\LoesungLuecke{return x;}{10cm}@*)

    (*@\LoesungLuecke{return x;}{10cm}@*)
    
    public double getX() {
    
        (*@\LoesungLuecke{return x;}{10cm}@*)
    }
    public double getY() {
    
        (*@\LoesungLuecke{return y;}{10cm}@*)
    }
    public void run() {
    
        (*@\LoesungLuecke{x += speedX;}{10cm}@*)
        
        (*@\LoesungLuecke{y += speedY;}{10cm}@*)
        
        (*@\LoesungLuecke{y += speedY;}{10cm}@*)
    }
    
    public void (*@\LoesungLuecke{setX(double wert}{10cm}@*) {
    
        (*@\LoesungLuecke{x = wert;}{10cm}@*)
    }
    
    public void (*@\LoesungLuecke{setY(double wert}{10cm}@*) {
    
        (*@\LoesungLuecke{y = wert;}{10cm}@*)
    }
    
    public String (*@\LoesungLuecke{getText(double wert}{10cm}@*) {
    
        (*@\LoesungLuecke{y = wert;}{15cm}@*)
        
        (*@\LoesungLuecke{y = wert;}{15cm}@*)
    }
    
    public void (*@\LoesungLuecke{setY(double wert}{10cm}@*) {
    
        (*@\LoesungLuecke{y = wert;}{15cm}@*)
        
        (*@\LoesungLuecke{y = wert;}{15cm}@*)
    }
}
\end{lstlisting}